\begin{definition}[An actual strip beside a flux endpoint]
    \label{def:peps_torus_charge_flux_crossing_strip}
    \lean{TNLean.PEPS.torusChargeFluxCrossingRegion,
        TNLean.PEPS.torusChargeFluxCrossingChargeBond}
    \leanok
    \uses{def:peps_three_plaquette_geometry,
        def:peps_torus_charge_string_edge}
    For $v=(x,y)$, take the translated four-column, two-row strip with
    lower-left vertex $(x-1,y+1)$. Its charge bond joins
    $(x+1,y+1)$ to $(x+2,y+1)$ on the lower row.
    Its leftmost plaquette contains the endpoint of the literal flux string.
    The tree consists of all horizontal bonds and the rightmost vertical bond.
    This is auxiliary geometry for the charge--flux figure of
    \cite{Schuch2010PEPS}, lines 2560--2581.
\end{definition}

\begin{theorem}[A physical charge crossing beside the literal flux endpoint]
    \label{thm:peps_torus_charge_flux_crossing}
    \lean{TNLean.PEPS.IsGIsometric.exists_unitary_torusChargeFluxCrossing,
        TNLean.PEPS.IsGIsometric.exists_unitary_torusCorrelatedChargeFluxCrossing}
    \leanok
    \uses{def:peps_torus_charge_flux_crossing_strip,
        thm:peps_regular_charge_flux_physical_transport,
        thm:peps_torus_swept_string_transport}
    Let $G$ be finite, let the periods satisfy $w\geq5$ and $H\geq4$, and
    let the native four-leg tensor be regular $G$-isometric. For each
    translated strip there is an original-spin unitary, chosen before the
    flux $k$, the scalar function $\chi$, the parameter $p$, and every
    crossing configuration. On the actual open contraction with the literal
    initial flux string, it changes the charge weight by
    \[
        \chi(p\eta_c)\longmapsto\chi(pk^{-1}\eta_c).
    \]
    The same construction is uniform in an unsummed partner reference $q$:
    \[
        \chi(pq^{-1}\eta_c)\longmapsto\chi(pq^{-1}k^{-1}\eta_c).
    \]
    Thus correlation is retained in its literal order. These statements
    include strips crossing either periodic seam and require no supplied
    tree gauge, Gram identity, or state hypothesis.

    This is an auxiliary physical operation beside an actual flux endpoint.
    It is not the identification of the complete prescribed braid or the
    reunion experiment of \cite{Schuch2010PEPS}, lines 2560--2615.
\end{theorem}
\begin{proof}\leanok
    The actual seven tree transports are the identity except for the
    rightmost upward transport $k$. The normalized tree gauge is therefore
    the identity on the lower row and $k$ on the upper row. The charge tail
    has trivial gauge. The leftmost cycle residual is $k^{-1}$ in the usual
    native order and $k$ when the upward bond reverses at the seam.
    Apply the original-spin reference-cycle unitary, using its adjoint
    in the first orientation. Uniformity in the coefficient function gives
    the correlated formula by substituting $t\mapsto\chi(pq^{-1}t)$.
\end{proof}
